% B's two content frames; F may insert them into the shared Beamer deck.
% Updated 2026-10-20 after F's actual 30-question S0 answers (PR #42).
% This is not an E scoring result or an S2 comparison.
\section{S0 普通检索基线}
\presenter{YaoYang-123456}
\source{B: reports/s0-local-run-20261019.md; F: PR \#42, s0-answer-summary-20261019.json}
\begin{frame}{S0：30题已完成真实检索与回答}
\small
\textbf{流程：}原始对话分块 $\rightarrow$ MiniLM 向量 $\rightarrow$
top-60 检索 $\rightarrow$ F 统一回答 $\rightarrow$ E 统一评分
\par\vspace{3mm}
\begin{tabularx}{\linewidth}{@{}p{33mm}Y@{}}
\toprule
固定数据 & 3个开发会话，30题（类别8/8/7/7）；前5题为联调子集\\[1mm]
B 的真实检索 & 384维、归一化、CPU；5/5和30/30题完成，每题60块；保留原始ID、时间和说话人\\[1mm]
F 的真实回答 & 5＋25次调用，30/30条生成成功；0次失败、截断或重复\\[1mm]
质量与验证 & A/B/F共29项针对性测试通过；E尚未判分，不能报告正确率\\\bottomrule
\end{tabularx}
\aside{S0是原始对话向量RAG，不是论文的平铺事实记忆消融；生成成功不等于答对。}
\note{建议约45秒。我们用普通向量检索作为StructMem的系统级基线。A冻结30题，B从原始会话检索，每题60块；来源、时间和说话人可追踪。F随后用统一模型按前5题、后25题完成30个真实回答，没有重复收费。但E还未评分，所以不能报告准确率。}
\end{frame}
